\documentclass[11pt,a4paper]{article}

% ============================================================
% Packages
% ============================================================
\usepackage[utf8]{inputenc}
\usepackage[T1]{fontenc}
\usepackage{lmodern}
\usepackage{amsmath,amssymb,amsthm,amsfonts,mathtools}
\usepackage{graphicx}
\usepackage{geometry}
\usepackage{booktabs}
\usepackage{tabularx}
\usepackage{caption}
\usepackage{subcaption}
\usepackage{setspace}
\usepackage{natbib}
\usepackage{hyperref}
\usepackage{xcolor}
\usepackage{microtype}
\usepackage{enumitem}

\geometry{margin=1in}
\onehalfspacing

\hypersetup{
    colorlinks=true,
    linkcolor=blue!60!black,
    citecolor=blue!60!black,
    urlcolor=blue!60!black
}

% ============================================================
% Theorem environments
% ============================================================
\newtheorem{theorem}{Theorem}
\newtheorem{proposition}{Proposition}
\newtheorem{lemma}{Lemma}
\newtheorem{corollary}{Corollary}
\theoremstyle{definition}
\newtheorem{definition}{Definition}
\theoremstyle{remark}
\newtheorem{remark}{Remark}

% ============================================================
% Mathematical notation
% ============================================================
\newcommand{\R}{\mathbb{R}}
\newcommand{\U}{\mathcal{U}}
\newcommand{\Ugrid}{\widehat{\mathcal{U}}}
\newcommand{\W}{\mathcal{W}}
\newcommand{\CVaR}{\operatorname{CVaR}}
\newcommand{\VaR}{\operatorname{VaR}}
\newcommand{\ESS}{\operatorname{ESS}}
\newcommand{\argmax}{\operatorname*{arg\,max}}
\newcommand{\argmin}{\operatorname*{arg\,min}}

% ============================================================
% Title
% ============================================================
\title{
State-robust portfolio selection with kernel-weighted conditional CVaR:
a semi-infinite model and grid-restricted exchange algorithm
}

\author{
Madani Bezoui$^{*,1}$ and Thiziri Sifaoui$^{2,3}$\\[0.5em]
\small $^1$CESI LINEACT, UR 7527, Nancy, France\\
\small $^2$Department of Mathematics and Computer Science,\\
\small University of Amine Elokkal El Hadj Moussa Eg Akhamouk,
Tamanghasset, Algeria\\
\small $^3$LAROMAD, Faculty of Sciences, UMMTO, Tizi Ouzou, Algeria
}

% Use a fixed submission date, not \today.
\date{August 15, 2026}

\begin{document}

\maketitle

\begin{NoHyper}
\def\thefootnote{*}
\footnotetext{
Corresponding author: Madani Bezoui,
\href{mailto:mbezoui@cesi.fr}{mbezoui@cesi.fr}.
}
\end{NoHyper}

% ============================================================
% Abstract
% ============================================================
\begin{abstract}
We study a state-robust portfolio selection problem in which observable
financial covariates index a family of kernel-weighted empirical return
distributions. The state vector consists of implied volatility and trailing
equity-market drawdown. For each state, a Nadaraya--Watson weighting scheme
defines a conditional empirical distribution, and portfolio risk is measured
by conditional value-at-risk (CVaR). The theoretical model minimizes the
largest conditional CVaR over a compact continuous state domain, resulting in
a convex semi-infinite program.

We establish continuity, convexity, and existence of an optimal solution.
For computation, we distinguish the continuous model from its finite-grid
approximation. The implemented algorithm is a constraint-generation method
over a prescribed state grid: it alternates between a finite linear programming
master problem and exhaustive separation over the grid. We prove finite
termination for the grid-restricted problem and provide a continuous-domain
error bound when conditional CVaR is uniformly Lipschitz in the state
variables.

The empirical analysis compares the proposed method with equal weighting,
minimum variance, unconditional CVaR, and a finite-regime CVaR benchmark in a
strictly chronological rolling-window design. All tables and figures are
generated directly from archived machine-readable outputs. The empirical
purpose is to evaluate whether continuous state indexing changes allocations
and realized risk, not to presume that state robustness must dominate simpler
benchmarks out of sample.
\end{abstract}

\noindent\textbf{Keywords:}
portfolio optimization; conditional value-at-risk; side information;
semi-infinite programming; kernel weighting; constraint generation.

\noindent\textbf{JEL classification:}
C44; C61; G11.

% ============================================================
\section{Introduction}
% ============================================================

Mean--variance portfolio optimization is sensitive to estimation errors,
especially errors in expected returns
\citep{markowitz1952portfolio,goldfarb2003robust}.
Robust and distributionally robust optimization address this problem by
protecting decisions against parameter or distributional uncertainty
\citep{delage2010distributionally,esfahani2018data}.
These formulations are often unconditional, although return distributions
and dependence structures vary with observable market conditions.

Conditional stochastic optimization incorporates side information into
data-driven decisions
\citep{bertsimas2020predictive,esteban2022distributionally,
nguyen2025robustifying}.
A common alternative in finance is to divide observations into a small
number of regimes \citep{ang2002international}. Such partitions are
interpretable but may introduce artificial discontinuities at regime
boundaries.

This paper studies an alternative construction. Observable market states
index a family of kernel-weighted empirical distributions. Instead of
optimizing only for the currently observed state, we seek one portfolio that
controls conditional CVaR across an entire state domain. We therefore call
the model \emph{state robust}. It is not a distributionally robust model:
conditional on a state, the empirical distribution is fixed by the selected
kernel and bandwidth.

The theoretical formulation is a convex semi-infinite program because the
constraint index set is continuous. The computational implementation,
however, uses a finite state grid. This distinction is important. Exhaustive
evaluation of a finite grid does not constitute certified continuous global
optimization unless a separate discretization-error bound is established.

Our contributions are:

\begin{enumerate}[label=(\roman*)]
    \item a kernel-weighted state-robust CVaR formulation with a compact
    portfolio set and compact state domain;

    \item continuity, convexity, and existence results for the continuous
    problem;

    \item a constraint-generation algorithm that terminates finitely for a
    fixed finite state grid;

    \item a uniform approximation bound connecting the grid problem to the
    continuous model under an explicit Lipschitz assumption; and

    \item a chronological empirical protocol comparing the method with
    transparent portfolio benchmarks.
\end{enumerate}

The term \emph{adaptive} refers exclusively to adaptive generation of active
state constraints. It does not mean that the final portfolio is an explicit
decision rule depending on the contemporaneous state.

% ============================================================
\section{Related literature}
% ============================================================

The paper intersects robust portfolio optimization, conditional stochastic
optimization, and semi-infinite programming.

Robust portfolio models protect allocations against parameter uncertainty
\citep{goldfarb2003robust,fabozzi2007robust}. CVaR provides a coherent and
computationally tractable tail-risk measure
\citep{rockafellar2000optimization,rockafellar2002conditional}.
Distributionally robust models extend this perspective to ambiguity sets of
probability distributions
\citep{delage2010distributionally,esfahani2018data}.

Side-information methods construct decisions using conditional empirical
distributions or predictive weights
\citep{bertsimas2020predictive,scaillet2004nonparametric}.
Recent conditional-DRO formulations additionally protect the conditional
distribution against misspecification
\citep{esteban2022distributionally,nguyen2025robustifying}.
Our formulation differs because it fixes one kernel-weighted empirical
distribution for each state and protects against variation over the state
index itself.

Semi-infinite programs contain finitely many decision variables and
infinitely many constraints
\citep{hettich1993semi,lopez2007semi}.
Exchange and cutting-plane algorithms iteratively add violated constraints.
Exact convergence statements depend critically on the quality of the
separation oracle. Inexact-oracle results require explicit error assumptions
and cannot be inferred solely from the use of a numerical grid
\citep{oustry2025convex}.

% ============================================================
\section{State-dependent empirical distributions}
% ============================================================

\subsection{Returns and lagged states}

Let

\[
x_t=(x_{1t},\ldots,x_{Nt})^\top\in\R^N
\]
denote the asset-return vector from trading day \(t-1\) to day \(t\).
Let

\[
y_{t-1}=(\log V_{t-1},D_{t-1})^\top\in\R^2
\]
be the information available before \(x_t\) is realized, where \(V_t\)
denotes VIX.

We define positive trailing drawdown by
\begin{equation}
D_t
=
1-
\frac{P_t}
{\max_{s\in\{\max(1,t-L),\ldots,t\}}P_s},
\qquad L=63,
\label{eq:drawdown}
\end{equation}
so that \(D_t\in[0,1]\). The market index \(P_t\), its source, and its
construction must be fixed before the backtest and documented in the
replication files.

The lagged pairing \((y_{t-1},x_t)\) prevents the use of information observed
after the return interval begins.

\subsection{State domain}

For a given training sample, define
\begin{equation}
\U
=
[\underline v,\overline v]
\times
[0,\overline d],
\label{eq:state-domain}
\end{equation}
where

\[
\underline v
=
v_{\min}-\delta_v,\qquad
\overline v
=
v_{\max}+\delta_v,
\]
and

\[
\overline d
=
\min\{1,d_{\max}+\delta_d\}.
\]

The bounds must be calculated using the current training window only. Using
full-sample extrema would introduce look-ahead bias.

A rectangular domain may include weakly supported state combinations.
Accordingly, we report \(\ESS(\theta)\) and conduct sensitivity analysis with
a supported domain
\begin{equation}
\U_E
=
\{\theta\in\U:\ESS(\theta)\ge E_{\min}\}.
\label{eq:supported-domain}
\end{equation}

\subsection{Kernel weights}

Let \(H\in\R^{2\times2}\) be symmetric positive definite. For
\(\theta\in\U\), define
\begin{equation}
p_t(\theta)
=
\frac{
\exp\!\left[
-\frac12(y_{t-1}-\theta)^\top H^{-1}(y_{t-1}-\theta)
\right]
}{
\sum_{s=1}^{T}
\exp\!\left[
-\frac12(y_{s-1}-\theta)^\top H^{-1}(y_{s-1}-\theta)
\right]
}.
\label{eq:kernel-weights}
\end{equation}

The common normalizing factor of the Gaussian kernel cancels from the ratio.
For numerical stability, the implementation evaluates
\eqref{eq:kernel-weights} using a log-sum-exp calculation.

The baseline bandwidth is
\begin{equation}
H
=
T^{-2/(M+4)}\widehat\Sigma_y
=
T^{-1/3}\widehat\Sigma_y,
\qquad M=2.
\label{eq:bandwidth}
\end{equation}

Bandwidth sensitivity is evaluated through

\[
H_\gamma=\gamma H,
\qquad
\gamma\in\{0.5,0.75,1,1.5,2\}.
\]

The effective sample size is
\begin{equation}
\ESS(\theta)
=
\frac{1}{\sum_{t=1}^{T}p_t(\theta)^2}.
\label{eq:ess}
\end{equation}

Because CVaR uses a lower-tail probability \(\tau\), we also report

\[
\ESS_{\mathrm{tail}}(\theta)=\tau\ESS(\theta).
\]
A strictly positive Gaussian weight does not imply adequate local support.

% ============================================================
\section{State-robust CVaR model}
% ============================================================

\subsection{Feasible portfolios}

Let \(\bar w\in(0,1]\) denote the maximum asset weight. Define
\begin{equation}
\W
=
\left\{
w\in\R^N:
\mathbf1^\top w=1,\;
0\le w_i\le\bar w,\;
\widehat\mu^\top w\ge\mu_{\mathrm{target}}
\right\}.
\label{eq:portfolio-set}
\end{equation}

The condition \(N\bar w\ge1\) is necessary for the capped simplex to be
nonempty. Let

\[
\widehat\mu_{(1)}
\ge
\widehat\mu_{(2)}
\ge\cdots\ge
\widehat\mu_{(N)}
\]
be the ordered estimated asset means, and define

\[
q=\left\lfloor\frac1{\bar w}\right\rfloor,
\qquad
r=1-q\bar w.
\]
When \(r>0\),
\begin{equation}
\mu_{\max}
=
\bar w\sum_{i=1}^{q}\widehat\mu_{(i)}
+
r\widehat\mu_{(q+1)}.
\label{eq:max-feasible-return}
\end{equation}
When \(r=0\), the residual term is omitted. Feasibility requires
\(\mu_{\mathrm{target}}\le\mu_{\max}\).

\subsection{Conditional CVaR}

For portfolio \(w\), define losses

\[
\ell_t(w)=-x_t^\top w.
\]

At confidence level \(\alpha=1-\tau\), conditional CVaR is
\begin{equation}
\Phi_\tau(w,\theta)
=
\min_{z\in\R}
\left\{
z+
\frac1{\tau}
\sum_{t=1}^{T}
p_t(\theta)[\ell_t(w)-z]_+
\right\}.
\label{eq:conditional-cvar}
\end{equation}

Any minimizer satisfies
\begin{equation}
z^*(w,\theta)
\in
\VaR_{1-\tau}(\ell(w)\mid\theta).
\label{eq:var-level}
\end{equation}
For a discrete weighted distribution, the VaR minimizer need not be unique,
although the optimized CVaR value is well defined.

\subsection{Continuous state-robust problem}

Define
\begin{equation}
G(w)=\max_{\theta\in\U}\Phi_\tau(w,\theta).
\label{eq:robust-objective}
\end{equation}

The theoretical continuous model is
\begin{equation}
v^*
=
\min_{w\in\W}G(w).
\label{eq:continuous-problem}
\end{equation}

Equivalently,
\begin{align}
\min_{w,\eta}\quad
&\eta
\\
\text{s.t.}\quad
&\Phi_\tau(w,\theta)\le\eta,
&&\forall\theta\in\U,
\\
&w\in\W.
\label{eq:sip}
\end{align}

\begin{proposition}[Well-posedness]
\label{prop:wellposed}
Suppose that \(H\) is symmetric positive definite, \(\U\) is nonempty and
compact, and \(\W\) is nonempty. Then:
\begin{enumerate}[label=(\roman*)]
    \item \(p_t(\theta)\) is positive and continuously differentiable on
    \(\U\);
    \item \(\Phi_\tau\) is jointly continuous on \(\W\times\U\);
    \item \(\Phi_\tau(\cdot,\theta)\) is convex on \(\W\);
    \item \(G\) is continuous and convex on \(\W\); and
    \item problem \eqref{eq:continuous-problem} admits an optimal solution.
\end{enumerate}
\end{proposition}

\begin{proof}
The numerator and denominator in \eqref{eq:kernel-weights} are smooth and the
denominator is strictly positive. Hence \(p_t\) is smooth.

Because \(\W\) is compact, portfolio losses are uniformly bounded. The
minimization over \(z\) in \eqref{eq:conditional-cvar} can therefore be
restricted to a common compact interval containing all attainable losses.
Continuity follows from Berge's maximum theorem. Convexity in \(w\) follows
from partial minimization of a jointly convex function in \((w,z)\).

The maximum of a jointly continuous function over compact \(\U\) is
continuous in \(w\). It is convex because it is a pointwise maximum of convex
functions. Existence follows from continuity of \(G\) and compactness of
\(\W\).
\end{proof}

% ============================================================
\section{Grid-restricted exchange algorithm}
% ============================================================

\subsection{Finite approximation}

Let

\[
\Ugrid_K=\{\theta^1,\ldots,\theta^K\}\subset\U
\]
be a finite grid. The implemented optimization problem is
\begin{equation}
\widehat v_K
=
\min_{w\in\W}
\widehat G_K(w),
\qquad
\widehat G_K(w)
=
\max_{\theta\in\Ugrid_K}
\Phi_\tau(w,\theta).
\label{eq:grid-problem}
\end{equation}

Problem \eqref{eq:grid-problem} is a finite robust linear program after
introducing state-specific VaR and excess-loss variables.

\subsection{Master problem}

At iteration \(k\), let

\[
\U_k\subseteq\Ugrid_K
\]
be the active state set. The master problem is
\begin{align}
\mathrm{LB}_k
=
\min_{w,\eta,z,u}\quad
&\eta
\\
\text{s.t.}\quad
&z_\theta+
\frac1{\tau}
\sum_{t=1}^{T}p_t(\theta)u_{t\theta}
\le\eta,
&&\forall\theta\in\U_k,
\\
&u_{t\theta}\ge -x_t^\top w-z_\theta,
&&t=1,\ldots,T,\;\theta\in\U_k,
\\
&u_{t\theta}\ge0,
&&t=1,\ldots,T,\;\theta\in\U_k,
\\
&w\in\W.
\label{eq:master}
\end{align}

For the resulting portfolio \(w_k\), exhaustive grid separation computes
\begin{equation}
\widehat G_k
=
\max_{\theta\in\Ugrid_K}\Phi_\tau(w_k,\theta)
\label{eq:grid-separation}
\end{equation}
and selects a maximizer \(\widehat\theta_k\).

The stopping residual is
\begin{equation}
r_k=\widehat G_k-\mathrm{LB}_k.
\label{eq:grid-residual}
\end{equation}

It is a residual for the finite-grid problem, not automatically a certified
optimality gap for the continuous problem.

\begin{proposition}[Finite termination on a fixed grid]
\label{prop:grid-termination}
Suppose every master problem is solved globally and every state in
\(\Ugrid_K\) is evaluated exactly during separation. If a newly violated state
is added whenever \(r_k>\epsilon\), then the algorithm terminates after at
most \(K\) state additions. At termination,

\[
0\le
\widehat G_K(w_k)-\widehat v_K
\le\epsilon.
\]
\end{proposition}

\begin{proof}
Each nonterminating iteration adds a state not already represented in the
master problem. Because \(\Ugrid_K\) contains \(K\) states, at most \(K\)
distinct additions are possible. Moreover,

\[
\mathrm{LB}_k\le\widehat v_K
\le\widehat G_K(w_k).
\]
Thus, when
\(\widehat G_K(w_k)-\mathrm{LB}_k\le\epsilon\),

\[
\widehat G_K(w_k)-\widehat v_K
\le
\widehat G_K(w_k)-\mathrm{LB}_k
\le\epsilon.
\]
\end{proof}

\subsection{Continuous-domain approximation}

Define the grid dispersion radius
\begin{equation}
\rho_K
=
\max_{\theta\in\U}
\min_{\widehat\theta\in\Ugrid_K}
\|\theta-\widehat\theta\|.
\label{eq:dispersion}
\end{equation}

\begin{proposition}[Grid approximation bound]
\label{prop:grid-bound}
Suppose that, uniformly for \(w\in\W\),

\[
|\Phi_\tau(w,\theta)-\Phi_\tau(w,\theta')|
\le
L_\Phi\|\theta-\theta'\|
\]
for all \(\theta,\theta'\in\U\). Then
\begin{equation}
\widehat G_K(w)
\le
G(w)
\le
\widehat G_K(w)+L_\Phi\rho_K
\label{eq:uniform-grid-bound}
\end{equation}
for every \(w\in\W\), and
\begin{equation}
\widehat v_K
\le
v^*
\le
\widehat v_K+L_\Phi\rho_K.
\label{eq:value-bound}
\end{equation}
If the exchange algorithm terminates with residual at most \(\epsilon\), then
\begin{equation}
G(w_k)-v^*
\le
\epsilon+L_\Phi\rho_K.
\label{eq:candidate-bound}
\end{equation}
\end{proposition}

\begin{proof}
For each \(\theta\in\U\), choose a nearest grid point
\(\widehat\theta\in\Ugrid_K\). Lipschitz continuity gives

\[
\Phi_\tau(w,\theta)
\le
\Phi_\tau(w,\widehat\theta)+L_\Phi\rho_K
\le
\widehat G_K(w)+L_\Phi\rho_K.
\]
Taking the maximum over \(\theta\) establishes
\eqref{eq:uniform-grid-bound}. Minimizing over \(w\) gives
\eqref{eq:value-bound}. Finally,

\[
G(w_k)
\le
\widehat G_K(w_k)+L_\Phi\rho_K
\le
\widehat v_K+\epsilon+L_\Phi\rho_K
\le
v^*+\epsilon+L_\Phi\rho_K.
\]
\end{proof}

\begin{remark}
A numerical value of \(L_\Phi\rho_K\) is informative only if its units,
coordinate system, norm, and magnitude relative to portfolio CVaR are
reported. A very large bound is mathematically valid but practically
vacuous.
\end{remark}

% ============================================================
\section{Empirical protocol}
% ============================================================

\subsection{Data}

The empirical universe consists of the Kenneth French 30 value-weighted
industry portfolios. VIX is obtained from Cboe. The broad-market series used
to construct \eqref{eq:drawdown} must be identified exactly.

The final manuscript must report:

\begin{itemize}
    \item exact source URLs and access dates;
    \item whether the market series is CRSP, a French market factor, or an
    industry-portfolio proxy;
    \item whether any VIX--VXO splice is used and, if so, its exact
    construction;
    \item the first and final aligned dates;
    \item the number of observations; and
    \item a checksum for the frozen input dataset.
\end{itemize}

No observation after a rebalancing date may be used to determine state bounds,
bandwidths, target returns, or portfolio weights.

\subsection{Rolling design}

For rebalancing date \(q\), use a trailing training window of 1,260 trading
days. The portfolio is evaluated over the following nonoverlapping holding
period.

A machine-readable backtest calendar must contain:

\begin{center}
\begin{tabular}{llllrr}
\toprule
Window & Train start & Train end & Hold start & Hold end & Observations\\
\midrule
\multicolumn{6}{c}{Generated by the backtest script}\\
\bottomrule
\end{tabular}
\end{center}

Fixed 21-trading-day blocks must be described as 21-day holding periods, not
as exact calendar months.

\subsection{Benchmarks}

The benchmarks are:

\begin{enumerate}
    \item monthly rebalanced \(1/N\);
    \item capped minimum variance;
    \item capped unconditional empirical CVaR; and
    \item capped four-regime CVaR using training-sample median splits.
\end{enumerate}

All optimized models use the same feasible set \(\W\). Solver failure must not
silently remove the target-return constraint. Every optimization exports its
termination status and feasibility residual.

\subsection{Turnover and transaction costs}

Let \(g_{iq}\) be asset \(i\)'s gross growth during holding period \(q\). The
weights immediately before rebalance \(q+1\) are
\begin{equation}
w^-_{i,q+1}
=
\frac{w^+_{i,q}g_{iq}}
{\sum_{j=1}^{N}w^+_{j,q}g_{jq}}.
\label{eq:drifted-weights}
\end{equation}

Turnover is
\begin{equation}
\mathrm{TO}_{q+1}
=
\frac12
\sum_{i=1}^{N}
|w^+_{i,q+1}-w^-_{i,q+1}|.
\label{eq:turnover}
\end{equation}

The drift must use the period ending at the current rebalancing date, not the
future holding-period return.

For proportional one-way cost \(c\),
\begin{equation}
r^{\mathrm{net}}_{q+1}
=
r^{\mathrm{gross}}_{q+1}
-
c\,\mathrm{TO}_{q+1}.
\label{eq:net-return}
\end{equation}

\subsection{Inference}

Sharpe-ratio comparisons use paired circular moving-block resampling of the
two strategy return series. Unless an inner standard-error estimator is
computed in every replication, the procedure must not be called studentized.

The manuscript reports:

\begin{itemize}
    \item the block length;
    \item number of replications;
    \item random seed;
    \item confidence-interval construction;
    \item definition of the empirical \(p\)-value; and
    \item multiplicity-adjusted results for several benchmarks.
\end{itemize}

% ============================================================
\section{Empirical results}
% ============================================================

All results in this section must be generated directly by the archived
analysis pipeline. Numerical values must not be manually entered into plotting
scripts.

\subsection{Performance}

\IfFileExists{../generated/performance_table.tex}{
    \input{../generated/performance_table.tex}
}{
    \begin{center}
    \fbox{
    \parbox{0.9\textwidth}{
    \textbf{Required generated file:}
    \texttt{generated/performance\_table.tex}.\\
    This table must be regenerated after correcting the data, state
    construction, grid definition, turnover timing, and solver-status
    handling.
    }}
    \end{center}
}

The interpretation must distinguish absolute return from downside protection.
If the proposed method has lower Sharpe, higher CVaR, or deeper drawdown than
a benchmark, this must be stated directly.

\subsection{Crisis-period results}

\IfFileExists{../generated/crisis_table.tex}{
    \input{../generated/crisis_table.tex}
}{
    \begin{center}
    \fbox{
    \parbox{0.9\textwidth}{
    \textbf{Required generated file:}
    \texttt{generated/crisis\_table.tex}.\\
    Crisis dates must be fixed before inspecting strategy performance.
    Subperiod drawdown calculations must include initial wealth \(W_0=1\).
    }}
    \end{center}
}

\subsection{Grid resolution and computational diagnostics}

The minimum required resolutions are

\[
21\times21,\quad
41\times41,\quad
81\times81,\quad
161\times161.
\]

For each resolution, report:

\begin{itemize}
    \item grid objective value;
    \item exchange objective value;
    \item grid residual;
    \item solve time;
    \item number of master iterations;
    \item number of active states;
    \item portfolio \(L_1\) distance; and
    \item objective-value difference.
\end{itemize}

A large portfolio \(L_1\) distance must not be described as negligible merely
because objective values are close.

\subsection{ESS and bandwidth sensitivity}

Report the minimum, 5th percentile, median, mean, and maximum ESS of all active
states. Also report \(\tau\ESS\). Sensitivity analysis must include both
bandwidth scaling and minimum-ESS restrictions.

\subsection{Statistical inference}

\IfFileExists{../generated/bootstrap_table.tex}{
    \input{../generated/bootstrap_table.tex}
}{
    \begin{center}
    \fbox{
    \parbox{0.9\textwidth}{
    \textbf{Required generated file:}
    \texttt{generated/bootstrap\_table.tex}.\\
    The table must identify whether intervals are percentile, basic,
    studentized, or normal-approximation block-bootstrap intervals.
    }}
    \end{center}
}

Failure to reject equality must not be interpreted as evidence that strategies
are equivalent. It indicates only insufficient evidence of a difference under
the selected test.

% ============================================================
\section{Discussion}
% ============================================================

The continuous formulation offers an interpretable method for studying
portfolio risk over a family of covariate-indexed empirical distributions.
Its main limitation is statistical support in extreme state regions. A
Gaussian kernel assigns positive mass everywhere, but the effective number of
observations may nevertheless be very small.

The grid-restricted exchange algorithm can reduce the size of the master LP
relative to a monolithic finite-grid formulation. Its computational advantage
does not remove the curse of dimensionality: Cartesian grid size grows
exponentially with state dimension, and local kernel support deteriorates in
high dimensions.

The model should also not be expected to dominate unconditional CVaR in every
crisis. Worst-state conditional CVaR and realized drawdown are different
criteria. Empirical conclusions must therefore be based on the corrected
out-of-sample evidence rather than on the model's robust interpretation.

Potential extensions include data-supported state domains, conditional
ambiguity sets, turnover regularization, and certified global separation for
low-dimensional state spaces.

% ============================================================
\section{Conclusion}
% ============================================================

We formulated a state-robust portfolio problem that minimizes the maximum
kernel-weighted conditional CVaR over a compact market-state domain. The
continuous model is a convex semi-infinite program and admits an optimal
solution under standard compactness and bandwidth assumptions.

The practical implementation is a grid-restricted constraint-generation
algorithm. It terminates finitely for a fixed grid and yields an
\(\epsilon\)-optimal solution to the finite-grid problem. Approximation of the
continuous model additionally depends on the grid dispersion and a valid
uniform Lipschitz bound.

The empirical contribution must be evaluated after the complete pipeline has
been rerun with synchronized data, code, tables, and figures. In particular,
claims of superior downside protection or continuous global optimality should
be made only if supported by the regenerated evidence and a quantitatively
useful approximation certificate.

% ============================================================
% Data and code availability
% ============================================================
\section*{Data and code availability}

The replication archive is available at:

\begin{center}
\url{https://github.com/MadBezoui/Robust-SIP-Portfolio}
\end{center}

The final submitted version should identify a permanent release tag and commit
hash. The archive must include:

\begin{itemize}
    \item source-data documentation and checksums;
    \item the exact backtest calendar;
    \item solver statuses for every window;
    \item machine-readable portfolio weights and returns;
    \item raw bootstrap replications;
    \item convergence histories;
    \item a Julia \texttt{Manifest.toml};
    \item pinned Python dependencies; and
    \item scripts that generate every table and figure without manual values.
\end{itemize}

\section*{Declaration of competing interests}

The authors declare that they have no known competing financial interests or
personal relationships that could have influenced the work reported in this
paper.

\section*{Use of generative artificial intelligence}

If generative artificial intelligence was used for language editing, software
assistance, or manuscript preparation, its role should be disclosed according
to the target journal's policy. The authors remain responsible for all
mathematical statements, references, code, data, and empirical results.

% Do not use \nocite{*}; cite only references actually used.
\bibliographystyle{plainnat}
\bibliography{references}

\end{document}
